\chapter{Estado del arte}

\section{Mantenimiento predictivo y detección de anomalías}

Esta sección desarrollará el contexto de mantenimiento predictivo, monitorización
de condición y detección de anomalías en maquinaria rotativa.

Aspectos a cubrir:

\begin{itemize}
  \item diferencias entre mantenimiento correctivo, preventivo y predictivo;
  \item papel de las señales de vibración en rodamientos y motores;
  \item fallos comunes: pista interior, pista exterior, bola y desbalanceo;
  \item diferencias entre detección de anomalías, diagnóstico y pronóstico.
\end{itemize}

\section{Modelos de detección de anomalías}

El MVP comparará primero modelos clásicos y reproducibles:

\begin{itemize}
  \item Isolation Forest;
  \item One-Class SVM;
  \item Local Outlier Factor;
  \item PCA con error de reconstrucción.
\end{itemize}

Los autoencoders densos y modelos recurrentes se considerarán en fases
posteriores, cuando el pipeline de tensores esté validado.

En la extension run-to-failure, la lectura de estos modelos cambia. Un detector
no se evalua solo por separar ventanas normales y anomalias, sino por producir
una trayectoria de score que anticipe degradacion, mantenga falsas alarmas
nominales controladas y permita construir un indicador de salud interpretable.
Por ello, PCA con error de reconstruccion, Isolation Forest, One-Class SVM y
autoencoders densos se tratan como familias comparables dentro de un mismo
protocolo temporal, no como modelos aislados optimizados por F1.

\section{Run-to-failure, indicadores de salud y modelos de reconstruccion}

Los datasets run-to-failure representan una trayectoria historica de un activo
hasta fallo. En este tipo de problema conviene distinguir cuatro conceptos:
anomalia puntual, degradacion sostenida, indicador de salud y pronostico de
vida util restante. Una alerta aislada puede deberse a ruido o transitorios; la
degradacion defendible requiere persistencia temporal, tendencia del score y
coherencia con el tramo nominal.

Los Health Indicators son habituales en mantenimiento predictivo porque
condensan la evolucion del activo en una curva de salud. En un sistema
experimental local, este indicador debe declararse como evidencia auxiliar y no
como prediccion industrial cerrada. Su calidad puede evaluarse mediante caida
entre tramo inicial y final, monotonicidad, robustez frente al ruido, volatilidad
nominal y fuerza de tendencia durante la degradacion.

Los autoencoders densos aportan una extension natural para este escenario,
porque aprenden a reconstruir normalidad y convierten el error de
reconstruccion en un score de anomalia. Sin embargo, su uso exige suficientes
ventanas nominales de entrenamiento, control de fuga temporal, arquitectura
cerrada, evaluacion comparable y guardarrailes que impidan presentar un modelo
mas complejo como mejora automatica. En este TFM se incorporan como familia
PyTorch reproducible sobre features tabulares, dejando los autoencoders
recurrentes y el RUL completo como lineas futuras.

\section{Sistemas multiagente aplicados a pipelines de datos}

Esta sección describirá el uso de arquitecturas multiagente para dividir
responsabilidades entre nodos especializados. El enfoque del proyecto separa:

\begin{itemize}
  \item decisiones de alto nivel tomadas por agentes;
  \item validación de contratos mediante esquemas estrictos;
  \item ejecución determinista de transformaciones sobre datos;
  \item auditoría de rutas, configuraciones, métricas y artefactos.
\end{itemize}

\section{Datasets públicos de referencia}

El dataset principal será CWRU Bearing Dataset
\cite{cwru_bearing_data_center,cwru_normal_baseline,cwru_12k_drive_end}. NASA
IMS Bearing Dataset se reserva como extensión para escenarios de degradación
temporal \cite{nasa_pcoe_repository,nasa_ims_bearings_zip}.
